% Requires \usepackage{amsmath}. Citations use \citep{} (natbib/biblatex);
% replace the citation keys with those in your .bib file.

\subsection{Species distribution models}

We related cockle occurrence and biomass to environmental conditions, sampling
gear, and larval connectivity using spatial generalized linear mixed models
fitted with the \texttt{sdmTMB} package \citep{anderson2024} in \textsf{R}
\citep{rcore}. For each survey location $i$ we modelled two responses: the
presence--absence indicator $z_i\in\{0,1\}$ and the biomass density
$y_i\ge 0$ (g\,m$^{-2}$).

Presence was modelled with a Bernoulli distribution and a logit link,
\begin{equation}
z_i \sim \mathrm{Bernoulli}(\pi_i),
\qquad
\operatorname{logit}(\pi_i) = \eta_i ,
\end{equation}
and biomass with a Tweedie distribution (compound Poisson--gamma, power
parameter $1<p<2$) and a log link, which accommodates the zero-inflated,
right-skewed nature of biomass data,
\begin{equation}
y_i \sim \operatorname{Tweedie}(\mu_i,\phi,p),
\qquad
\log(\mu_i) = \eta_i ,
\end{equation}
where $\mu_i$ is the mean, $\phi$ the dispersion parameter, and $p$ the power
parameter.

Both responses shared the same linear predictor,
\begin{equation}
\eta_i = \beta_0 + \gamma_{g(i)}
      + \sum_{k=1}^{5}\beta_k\,x_{ik}
      + \theta\,c_i
      + \omega_{s_i},
\label{eq:linpred}
\end{equation}
where $\beta_0$ is the intercept; $\gamma_{g(i)}$ is the effect of the sampling
gear $g(i)$ (a categorical nuisance term); $x_{ik}$ are five environmental
covariates (depth, temperature, salinity, oxygen deficiency, and maximum bed
shear stress) with coefficients $\beta_k$; $c_i$ is a connectivity covariate
with coefficient $\theta$; and $\omega_{s_i}$ is a spatial random field
evaluated at the location $s_i$ of observation $i$. All continuous covariates
were standardized to zero mean and unit variance before fitting, so that the
estimated coefficients are directly comparable on the link scale.

\paragraph{Spatial random field.}
The spatial field $\boldsymbol{\omega}=(\omega_1,\dots)^\top$ was assumed to be a
zero-mean Gaussian random field with a Mat\'ern covariance function,
approximated by a Gaussian Markov random field through the stochastic partial
differential equation (SPDE) approach \citep{lindgren2011}. The field is
characterized by a marginal standard deviation $\sigma_\omega$ and a range
$\rho$ (the distance at which spatial correlation falls to approximately
$0.13$). Because the study area is a topologically complex fjord, we used a
barrier formulation \citep{bakka2019}: the Mat\'ern range was reduced to a
fraction (0.1) of $\rho$ over land, so that spatial correlation follows water
rather than crossing land barriers. The SPDE mesh was built with a minimum edge
length of 1\,km, and predicted land cells were fixed to zero biomass.

\paragraph{Connectivity covariates.}
Larval connectivity was derived from a biophysical particle-tracking model of
cockle larval dispersal \citep{yourABMref}. Let $F_{jk}$ denote the entries of
the resulting flow matrix, i.e.\ the probability that a larva released in cell
$j$ settles in cell $k$. From $F$ we computed, for each cell $k$, structural
connectivity metrics that depend only on the dispersal network and are therefore
independent of the cockle observations. The primary metric was the weighted
in-strength,
\begin{equation}
c_k = \sum_{j} F_{jk},
\end{equation}
the total incoming settlement probability. We additionally considered the
eigenvector centrality of $F$ and the log-transformed, biomass-weighted larval
supply $\log\!\big(1+\sum_j b_j F_{jk}\big)$, where $b_j$ is the predicted mean
biomass of source cell $j$. Model support was insensitive to the choice of
metric, and in-strength was the best supported; it is used in the final model.

\paragraph{Model selection.}
We compared six model structures obtained by crossing the presence or absence of
(i) the spatial field, (ii) the environmental covariates, and (iii) the
connectivity covariate, each retaining the gear term, and fitted separately for
presence and biomass. Models were fitted by maximum likelihood via Template
Model Builder \citep{kristensen2016} and ranked by Akaike's Information
Criterion (AIC); all models were fitted to the same complete-case data set so
that AIC values are comparable. Convergence was verified from the Hessian and
gradient checks reported by \texttt{sdmTMB}.

\paragraph{Variance partitioning.}
We summarized the relative contributions of the fixed effects and the spatial
field with the marginal and conditional coefficient of determination
\citep{nakagawa2017},
\begin{equation}
R^2_{\text{marginal}}
  = \frac{\sigma^2_f}{\sigma^2_f+\sigma^2_\omega+\sigma^2_d},
\qquad
R^2_{\text{conditional}}
  = \frac{\sigma^2_f+\sigma^2_\omega}{\sigma^2_f+\sigma^2_\omega+\sigma^2_d},
\end{equation}
where $\sigma^2_f=\operatorname{Var}\!\big(\hat{\beta}_0+\hat{\gamma}_{g(i)}
+\sum_k\hat{\beta}_k x_{ik}+\hat{\theta}c_i\big)$ is the variance of the
fixed-effects linear predictor, $\sigma^2_\omega$ is the variance of the spatial
field, and $\sigma^2_d$ is the distribution-specific (observation-level)
variance on the link scale. For the Tweedie model the latter was approximated by
$\sigma^2_d=\tfrac{1}{n}\sum_i\log\!\big(1+\phi\,\mu_i^{\,p-2}\big)$
\citep{nakagawa2017}. The fixed share $R^2_{\text{marginal}}$ was further divided
among the environmental, gear, and connectivity blocks in proportion to each
block's covariance with the fixed-effects linear predictor, so that the blocks
sum to $R^2_{\text{marginal}}$.

All analyses were carried out in \textsf{R} \citep{rcore}; models were fitted
with \texttt{sdmTMB} \citep{anderson2024} and barrier meshes constructed with
\texttt{sdmTMBextra} and \texttt{fmesher}.
